\documentclass[10pt,a4paper]{amsart}
\usepackage[margin=19mm]{geometry}
\usepackage{amsmath,amssymb,mathtools}
\usepackage[scr=rsfs,cal=euler]{mathalfa}
\usepackage{xfrac}
\usepackage[unicode,hidelinks,bookmarks=false]{hyperref}
\newcommand{\ie}{i.e.\ }
\newcommand{\cf}{cf.\ }
\newcommand{\resp}{respectively\ }
\newcommand{\Subsection}{\subsection}
\providecommand{\nobreakdash}{\nobreak\hbox{-}}
\setlength{\emergencystretch}{2em}
\multlinegap0pt
\usepackage{tikz}
\usetikzlibrary{cd}
\usepackage{mathtools}
\usepackage{amssymb}
\newtheorem{theorem}{Theorem}
\newcommand{\C}{\mathcal C}
\def\P{\mathcal P}
\def\Pcoalg{\P\mathsf{CoAlg}}
\newcommand{\V}{\mathcal V}
\newcommand{\eq}{\operatorname{eq}}
\def\id{\operatorname{id}}
\newcommand{\uC}{\underline{\C}}
\newcommand{\uPcoalg}{\underline{\Pcoalg}}
\newcommand{\uV}{\underline{\mathcal V}}
\def\unit{I}
\title{Enriched coalgebras are sometimes comonadic}
\author{Ois\'in Flynn-Connolly}
\date{Source: arXiv:2604.09354v3 (CC BY 4.0)}
\begin{document}
\maketitle

\noindent\emph{Source and licence.} Enriched coalgebras are sometimes comonadic. Version arXiv:2604.09354v3 (CC BY 4.0), distributed by arXiv under the Creative Commons Attribution 4.0 International licence, http://creativecommons.org/licenses/by/4.0/, as recorded in the arXiv licence metadata for this exact version and shown on the abstract page https://arxiv.org/abs/2604.09354v3. Full licence notice: \texttt{licenses/arxiv-2604.09354-CC-BY-4.0.txt}.\par\smallskip

\noindent\textit{Editorial scope.} Counted unit: the article's theorem enriefcat (source0.tex lines 658--661). The article's complete original proof of it is delivered with it (source0.tex lines 662--748, one \texttt{proof} environment of 87 lines). No mathematics is written, repaired or re-derived: the statement and the proof are the article's own lines, in the article's own notation, and no retained line is substituted or re-flowed.  No further source-local context is needed: every cross-reference of every retained line resolves inside the two delivered spans.  Nothing was omitted from any retained span, so the package carries no omission record.  No retained line contains a citation, so the wrapper carries no bibliography and no bibliography entry.  The retained lines contain the article's own diagram code, which the wrapper typesets with the standard tikz packages; no image file is delivered and the package declares no asset dependency.  Editorial formatting only, in addition: an AMS \texttt{amsart} preamble that restates the article's own declarations and macros used by the retained lines, namely the environments \texttt{theorem} on separate counters, together with the standard packages those lines need.  The retained environments print in this document as Theorem 1 in order of appearance, whereas the article prints the same items with the numbering of its own full document.  The complete native source of the article as distributed in the arXiv e-print for this exact version is bundled unchanged as \texttt{sources/198203/source0.tex}.
\begin{theorem}
\label{enriefcat}
    Let $\V$ be Cartesian and $\P$ be an operad (not necessarily unital). Then, with definitions as above, $\Pcoalg$ is a $\V$-category.
\end{theorem}

\begin{proof}
The base category $\V$ is complete, which ensure that the equaliser defining the mapping objects $\uPcoalg(X, Y)$ exists. To establish that $\Pcoalg$ forms a $\V$-category, we must show that these subobjects admit identity and composition morphisms. 

For any $\P$-coalgebra $X$, it is easy to see that the identity morphism $j_X \colon \unit_\V \to \uC(X, X)$ is strictly equalised by the maps $(f\circ \pi_r)^c$ and $(g\circ \pi_r)^c$.

For composition, we define the restricted map $c'$ as the composite in $\V$:
$$ c' \colon \uPcoalg(Y, Z) \times \uPcoalg(X, Y) \xrightarrow{\eq_{Y,Z} \times \eq_{X,Y}} \uC(Y, Z) \times \uC(X, Y) \xrightarrow{c} \uC(X, Z). $$
To show $c'$ factors uniquely to a composition $c_{\P}$, we must prove it equalises the diagram
$$
\begin{tikzcd}
\uC(X, Z) \ar[r,shift left=.75ex,"f"]
  \ar[r,shift right=.75ex,swap,"g"]
&
\prod_{r\geq 0} \uV(\P(r), \uC(X, Z^{\otimes r}) ) 
\end{tikzcd}
$$
defining $\uPcoalg(X, Z)$. It is again sufficient to check this property on adjoints of the equalisers. The trick is now to use the universal property of the equaliser to move the maps to $Y$.

First, we have the obvious diagram:

$$
\begin{tikzcd}[column sep=huge, row sep=huge]
\uPcoalg(Y, Z) \times \uPcoalg(X, Y) \times \P(r)
    \arrow[r, "\eq_{Y,Z} \times \eq_{X,Y} \times \id"]
    \arrow[d, "\id \times \id \times \phi_Z"']
& \uC(Y, Z) \times \uC(X, Y) \times \P(r)
    \arrow[d, "c \times \id"]
\\
\uPcoalg(Y, Z) \times \uPcoalg(X, Y) \times \uC(Z, Z^{\otimes r})
     \arrow[d, "\eq_{Y,Z} \times \eq_{X,Y} \times \id"]
& \uC(X, Z) \times \P(r)
    \arrow[d, "(g \circ \pi_r)^c"] 
\\
\uC(Y, Z) \times \uC(X, Y) \times \uC(Z, Z^{\otimes r}) \arrow[r, "c \circ (\id\times c)\circ s"]
& \uC(X, Z^{\otimes r})
\end{tikzcd}
$$
The map  $s$ just rearranges factors in the obvious way. 


By the defining property of the equaliser of $\eq_{Y,Z}$ we can replace the down-left path with:

$$
\begin{tikzcd}[column sep=huge, row sep=huge]
\uPcoalg(Y, Z) \times \uPcoalg(X, Y) \times \P(r)
    \arrow[r, "\eq_{Y,Z} \times \eq_{X,Y} \times \id"]
    \arrow[d, "\id \times \id \times \phi_Y"']
& \uC(Y, Z) \times \uC(X, Y) \times \P(r)
    \arrow[d, "c \times \id"]
\\
\uPcoalg(Y, Z) \times \uPcoalg(X, Y) \times \uC(Y, Y^{\otimes r})
     \arrow[d, "(\eq_{Y,Z}^{\times r}\circ \Delta^{(r)}) \times \eq_{X,Y} \times \id"]
& \uC(X, Z) \times \P(r)
    \arrow[d, "(g \circ \pi_r)^c"] 
\\
\uC(Y, Z)^{\times r} \times \uC(X, Y) \times \uC(Y, Y^{\otimes r}) 
 \arrow[r, "c \circ (T \times c)\circ s"]
& \uC(X, Z^{\otimes r})
\end{tikzcd}
$$
On the bottom we have once again switched factors. Now, we use the defining property of the equaliser of $\eq_{X,Y}$ to obtain:
$$
\begin{tikzcd}[column sep=huge, row sep=huge]
\uPcoalg(Y, Z) \times \uPcoalg(X, Y) \times \P(r)
    \arrow[r, "\eq_{Y,Z} \times \eq_{X,Y} \times \id"]
    \arrow[d, "\id \times \id \times \Phi_X"']
& \uC(Y, Z) \times \uC(X, Y) \times \P(r)
    \arrow[d, "c \times \id"]
\\
\uPcoalg(Y, Z) \times \uPcoalg(X, Y) \times \uC(X, X^{\otimes r})
     \arrow[d, "(\eq_{Y,Z}^{\times r}\circ \Delta^{(r)}) \times (\eq_{X,Y}^{\times r}\circ \Delta^{(r)} ) \times \id"]
& \uC(X, Z) \times \P(r)
    \arrow[d, "(g \circ \pi_r)^c"] 
\\
\uC(Y, Z)^{\times r} \times \uC(X, Y)^{\times r} \times \uC(X, X^{\otimes r}) 
 \arrow[r, "c \circ (T \times c)\circ s"]
& \uC(X, Z^{\otimes r})
\end{tikzcd}
$$



But now the left down path is precisely equal to $(f \circ \pi_r)^c \circ (c'\times \id)$. Thus, the restricted composition $c'$ strictly equalises the diagram as required.


The associativity and unitality axioms are directly inherited from the underlying $\V$-category $\C$, since composition is a restriction.    
\end{proof}

\end{document}
